\section*{第 \ref{ch:g10:common-ancestry} 章——身体构型与共同祖先}

\begin{solution}{exo:g10:common-ancestry:1}
两侧对称，具有前端和后端；具内骨骼，
内骨骼含脊柱，其上为背侧神经索；腹侧
为消化管，心脏位于其下方靠近前端处
（并有成对的肢或鳍）。
\end{solution}

\begin{solution}{exo:g10:common-ancestry:2}
结构和在\omterm{def:g10:common-ancestry:bodyplan}{身体构型}中的位置相同，而功能可以不同。四足动物的前肢：
肱骨、桡骨与尺骨、腕骨、指骨——
见于手臂、翅膀或鳍中。
\end{solution}

\begin{solution}{exo:g10:common-ancestry:3}
属于同功：功能相同（飞行），结构却完全不同——
鸟的是骨骼、肌肉和皮肤，蝴蝶的是外骨骼膜；
二者并不说明亲缘关系。
\end{solution}

\begin{solution}{exo:g10:common-ancestry:4}
蜥蜴：包含青蛙和蜥蜴的最小框是“四足动物”；
鳟鱼只在更大的“\omterm{def:g10:common-ancestry:bodyplan}{脊椎动物}”框内。
\end{solution}

\begin{solution}{exo:g10:common-ancestry:5}
具有中间性状组合的化石；\omterm{def:g10:common-ancestry:bodyplan}{脊椎动物}胚胎共有的发育阶段；
分子序列比较
与解剖学分组相吻合。
\end{solution}

\begin{solution}{exo:g10:common-ancestry:6}
哺乳动物（有毛发痕迹、乳汁；四足动物的肢；
用肺呼吸，且哺乳动物具有羊膜发育）。髋骨是后肢的残迹：
鲸类起源于四足的陆生哺乳动物。
\end{solution}

\begin{solution}{exo:g10:common-ancestry:7}
鳄鱼归入“羊膜动物”，与蜥蜴和鸽子同类（而不属于“哺乳动物”）；
蝾螈归入“四足动物”，在“羊膜动物”之外，
与青蛙同类。
\end{solution}

\begin{solution}{exo:g10:common-ancestry:8}
一个亲缘类群必须包含其共同祖先的全部后代；
鳟鱼和肺鱼的共同祖先，也是青蛙及所有四足动物的祖先；
若一个类群包含鳟鱼和肺鱼却不含青蛙，则遗漏了某些后代；
“鱼”指的是一种生活方式，
而非一个谱系。
\end{solution}

\begin{solution}{exo:g10:common-ancestry:9}
黑猩猩（0）、小鼠（9）、鸡（13）、青蛙（18）、金枪鱼（21）、酵母（45）：
依次为哺乳动物、羊膜动物、四足动物、\omterm{def:g10:common-ancestry:bodyplan}{脊椎动物}、真核生物——
即各框的顺序。
\end{solution}

\begin{solution}{exo:g10:common-ancestry:10}
成年鱼保留的结构会短暂出现在我们的胚胎中，
并被改造：鳃裂成为耳、咽喉和腺体的一部分，
尾部被吸收。胚胎经历某一祖先阶段，
正是遗传自共同祖先所预测的结果，
而为某种目的而设计则不会如此。
\end{solution}

\begin{solution}{exo:g10:common-ancestry:11}
若鸟类起源于类爬行动物的祖先，则某个祖先必须同时具有
爬行类特征和最初的鸟类特征；在两类群之间找到年代相符的
此类化石，即证实了这一预测。若无共同祖先，
兼具两类不同生物特征的生物便
没有存在的理由。
\end{solution}

\begin{solution}{exo:g10:common-ancestry:12}
属于同功：功能相同，总体设计相似，
但胚胎起源不同，结构的排列方式也相反。若属于\omterm{def:g10:common-ancestry:homology}{同源}，
则软体动物与\omterm{def:g10:common-ancestry:bodyplan}{脊椎动物}的共同祖先已具有照相机式眼；
若属于同功，则其祖先至多只有简单的感光斑块，
两个谱系各自独立地形成了照相机式眼。
\end{solution}

\begin{solution}{exo:g10:common-ancestry:13}
人与小鼠的共同祖先比它们各自与鸡的共同祖先更近；
两个谱系与鸡的谱系分开的时间相同，
因此相对于鸡，二者积累的变化数量大致相同。
人与鸡、小鼠与鸡之间的距离
度量的是同一段时间。
\end{solution}

\begin{solution}{exo:g10:common-ancestry:14}
蛇具有羊膜卵，以及蜥蜴式的骨骼和颅骨；
部分蛇类（如蟒类）仍保留后肢和骨盆的残迹；
也存在带有小腿的化石蛇。“四足动物”指
由四肢的祖先演化而来，无论其是否保留四肢。
\end{solution}

\begin{solution}{exo:g10:common-ancestry:15}
生活方式可以解释鳍的形状，
却无法解释鲸的鳍中为何含有与蝙蝠相同排列的手臂骨骼——
一块扁平的桨并不需要腕骨。它也无法解释人类胚胎中的鳃裂，
而人类胚胎对此并无用处。它同样无法解释，
一种对环境不可见的呼吸蛋白，为何在不同\omterm{def:g10:biodiversity-scales:species}{物种}之间的差异，
恰好与骨骼所预测的量相符。
只有遗传自共同祖先，才能解释这三者。
\end{solution}

\begin{solution}{pb:g10:common-ancestry:1}
\textbf{1.} 脊柱：\omterm{def:g10:common-ancestry:bodyplan}{脊椎动物}。

\textbf{2.} 四肢定义了青蛙、蜥蜴、鸽子、小鼠：
四足动物。羊膜卵定义了蜥蜴、鸽子、小鼠：
羊膜动物。

\textbf{3.} 不能。某个性状若只见于一个\omterm{def:g10:biodiversity-scales:species}{物种}，便不能划分出类群，它只能标识该\omterm{def:g10:biodiversity-scales:species}{物种}，
（若再加入更多亲缘\omterm{def:g10:biodiversity-scales:species}{物种}，
——其他鸟类、其他哺乳动物——则可界定一个类群）。

\textbf{4.} \omterm{def:g10:common-ancestry:bodyplan}{脊椎动物} $\supset$ 四足动物 $\supset$ 羊膜动物，
鳟鱼位于第二个方框之外，青蛙位于第三个方框之外；
蜥蜴、鸽子和小鼠位于第三个方框之内。

\textbf{5.} 青蛙的。骨骼相同，排列顺序相同。
它对应于“四肢”这一性状，而鱼鳍缺少该性状。

\textbf{6.} 蜥蜴和鸽子（12）。鳟鱼（与所有\omterm{def:g10:biodiversity-scales:species}{物种}的距离为 19--21）。

\textbf{7.} 鸽子（13）、蜥蜴（14）、青蛙（18）、鳟鱼（21）。

\textbf{8.} 蜥蜴（16）、鸽子（17）、小鼠（18）、鳟鱼（19）：
羊膜动物与鳟鱼的距离大致相等，鳟鱼则更远——这与方框的划分一致，
（青蛙与每一种羊膜动物的距离都比与鳟鱼的距离近）。

\textbf{9.} 蜥蜴和鸽子在各自谱系与小鼠的谱系分开之后，才彼此分开：
在羊膜动物之内，蜥蜴和鸽子的亲缘关系更近。
该表无法体现这一点：
表中没有任何一个性状仅为蜥蜴和鸽子所共有。

\textbf{10.} 鳟鱼的谱系在其余四个\omterm{def:g10:biodiversity-scales:species}{物种}的共同谱系中分出，
早于这四个\omterm{def:g10:biodiversity-scales:species}{物种}彼此分开：
每个\omterm{def:g10:biodiversity-scales:species}{物种}与鳟鱼分开的时间相同，
因此各自积累的差异数大致相同。

\textbf{11.} 鳟鱼位于其他所有\omterm{def:g10:biodiversity-scales:species}{物种}之外；四种四足动物聚为一群；
三种羊膜动物又在其中聚为一群。

\textbf{12.} 解剖结构与呼吸蛋白本无理由相互吻合，
除非二者记录的是同一段历史。两个独立的证据讲述同一个故事，
使偶然相似的解释变得难以成立。

\textbf{13.} 应将其归入“哺乳动物”，与小鼠并列（有毛发，有乳汁，且分子距离最小）。
其翅膀是由四足动物的前肢变化而来，与鸽子的翅膀属同功器官，
而非与有羽毛的飞行\omterm{def:g10:common-ancestry:homology}{同源}；
该分子（与鸽子距离 13，与小鼠距离 8）证实了这一点。

\textbf{14.} \omterm{def:g10:cells-common-unit:cell}{细胞}色素 c 变化过于缓慢，无法区分近缘\omterm{def:g10:biodiversity-scales:species}{物种}，
而且它并非造成可见差异的\omterm{def:g10:universal-dna:gene}{基因}。
它揭示的是深层的亲缘关系，
而非决定\omterm{def:g10:biodiversity-scales:species}{物种}特征的\omterm{def:g10:universal-dna:gene}{基因}。

\textbf{15.} 真核生物：具有\omterm{def:g10:cells-common-unit:organelle}{细胞核}和\omterm{def:g10:cells-common-unit:organelle}{线粒体}的\omterm{def:g10:cells-common-unit:cell}{细胞}，
它们共用这种呼吸蛋白。

\textbf{16.} 两个谱系在 3.2 亿年间共积累了 13 个差异，
即 $13/2 \approx 6.5$ 个／每条谱系，
即沿每条谱系约每 1 亿年积累 2 个。

\textbf{17.} 青蛙–小鼠为 18，青蛙–鸽子为 17：平均 17.5，即 8.75
个／每条谱系；按每 1 亿年 2 个计，约 4.3 亿年。

\textbf{18.} 鳟鱼的平均值约为 20：每条谱系 10 个，
约 5 亿年。

\textbf{19.} 估计值（4.3 亿、5 亿年）比化石年代（3.7 亿、4.2 亿以上）偏大，
但数量级正确，先后顺序也一致；
速率并非恒定不变，且计数经过取整，
因此两者在 20\% 以内相符，已是该方法所能达到的精度。

\textbf{20.} 二者确立了同样的嵌套亲缘关系——先是鳟鱼，再是青蛙，最后是三种羊膜动物——
而分子还提供了一把“时钟”：
可估计各分支分开至今已有多久。
\end{solution}
